\documentclass[11pt,twocolumn]{article}

\usepackage[utf8]{inputenc}
\usepackage[T1]{fontenc}
\usepackage{amsmath,amssymb}
\usepackage{graphicx}
\usepackage{booktabs}
\usepackage{hyperref}
\usepackage{listings}
\usepackage{xcolor}
\usepackage[margin=1in]{geometry}
\usepackage{fancyhdr}
\usepackage{abstract}

\definecolor{codegreen}{rgb}{0,0.6,0}
\definecolor{codegray}{rgb}{0.5,0.5,0.5}
\definecolor{codepurple}{rgb}{0.58,0,0.82}
\definecolor{backcolour}{rgb}{0.95,0.95,0.92}

\lstdefinestyle{mystyle}{
    backgroundcolor=\color{backcolour},
    commentstyle=\color{codegreen},
    keywordstyle=\color{magenta},
    numberstyle=\tiny\color{codegray},
    stringstyle=\color{codepurple},
    basicstyle=\ttfamily\footnotesize,
    breakatwhitespace=false,
    breaklines=true,
    captionpos=b,
    keepspaces=true,
    numbers=left,
    numbersep=5pt,
    showspaces=false,
    showstringspaces=false,
    showtabs=false,
    tabsize=2
}
\lstset{style=mystyle}

\title{\textbf{LuxPy: Python SDK for Data Science and ML}}
\subtitle{Pandas Integration, Jupyter Support}

\author{Zach Kelling\thanks{zach@hanzo.ai} \\ \textit{Hanzo Industries \quad Lux Industries \quad Zoo Labs Foundation} \\ \texttt{research@lux.network}}

\date{2024}

\begin{document}

\maketitle

\begin{abstract}
We present the Lux Python SDK, a client library designed for data science workflows, automated trading, and machine learning pipelines. The SDK provides native pandas integration for blockchain data analysis, async batch operations for high-throughput data collection, and seamless integration with popular ML frameworks. We introduce a novel streaming data architecture that enables real-time feature engineering from blockchain events. Benchmarks demonstrate processing of 100,000+ transactions per second for historical analysis and sub-second latency for live trading signals. The SDK bridges the gap between blockchain infrastructure and the Python data science ecosystem.
\end{abstract}

\section{Introduction}

Python dominates data science and machine learning, with pandas, NumPy, and scikit-learn forming the foundation of modern analytics workflows. Blockchain data analysis requires:

\begin{itemize}
    \item \textbf{Efficient data retrieval}: Batch fetching of historical transactions
    \item \textbf{Tabular representation}: DataFrames for exploratory analysis
    \item \textbf{Real-time streaming}: Live data for trading and monitoring
    \item \textbf{ML integration}: Feature engineering and model deployment
\end{itemize}

The Lux Python SDK addresses these requirements while maintaining idiomatic Python design. The SDK wraps the high-performance Rust core via PyO3, providing both Python convenience and native speed for computationally intensive operations.

This paper presents the SDK architecture, pandas integration patterns, and ML pipeline examples.

\section{Architecture}

\subsection{Package Structure}

The SDK follows Python packaging conventions:

\begin{lstlisting}[language=Python,caption=Package Structure]
luxsdk/
    __init__.py
    client.py       # High-level client
    wallet.py       # Key management
    chains/
        x_chain.py  # X-Chain operations
        p_chain.py  # P-Chain operations
        c_chain.py  # C-Chain (web3.py compat)
    data/
        frames.py   # DataFrame utilities
        streaming.py # Real-time data
        cache.py    # Local caching
    ml/
        features.py # Feature engineering
        signals.py  # Trading signals
    _core/          # Rust bindings (PyO3)
\end{lstlisting}

\subsection{Rust Core Integration}

Performance-critical operations use Rust via PyO3:

\begin{lstlisting}[language=Python,caption=Rust Bindings]
# Python wrapper around Rust implementation
from luxsdk._core import (
    sign_transaction,
    verify_signature,
    parse_utxos,
    serialize_tx,
)

class Signer:
    """Python interface to Rust signing."""

    def __init__(self, private_key: bytes):
        # Store key in Rust, never exposed to Python
        self._handle = _core.create_signer(private_key)

    def sign(self, message: bytes) -> bytes:
        """Sign message using Rust crypto."""
        return _core.sign(self._handle, message)

    def __del__(self):
        # Securely free Rust memory
        _core.free_signer(self._handle)
\end{lstlisting}

\subsection{Async Design}

The SDK uses \texttt{asyncio} for concurrent operations:

\begin{lstlisting}[language=Python,caption=Async Client]
import asyncio
from typing import AsyncIterator

class LuxClient:
    """Async client for Lux blockchain."""

    def __init__(
        self,
        endpoint: str,
        network: str = "mainnet",
    ):
        self.endpoint = endpoint
        self.network = network
        self._session: aiohttp.ClientSession | None = None

    async def __aenter__(self) -> "LuxClient":
        self._session = aiohttp.ClientSession()
        return self

    async def __aexit__(self, *args):
        if self._session:
            await self._session.close()

    async def get_balance(
        self,
        address: str,
        asset_id: str = "LUX",
    ) -> int:
        """Get balance for address."""
        result = await self._rpc_call(
            "avm.getBalance",
            {"address": address, "assetID": asset_id},
        )
        return int(result["balance"])

    async def stream_blocks(
        self,
        chain: str = "X",
    ) -> AsyncIterator[Block]:
        """Stream new blocks as they're produced."""
        async for event in self._subscribe(f"{chain}-chain/blocks"):
            yield Block.from_dict(event)
\end{lstlisting}

\section{Pandas Integration}

\subsection{DataFrame Conversion}

All data types provide DataFrame conversion:

\begin{lstlisting}[language=Python,caption=DataFrame Utilities]
import pandas as pd
from dataclasses import dataclass
from typing import List

@dataclass
class UTXO:
    tx_id: str
    output_index: int
    asset_id: str
    amount: int
    addresses: List[str]
    locktime: int

    @classmethod
    def to_dataframe(cls, utxos: List["UTXO"]) -> pd.DataFrame:
        """Convert UTXOs to DataFrame."""
        return pd.DataFrame([
            {
                "tx_id": u.tx_id,
                "output_index": u.output_index,
                "asset_id": u.asset_id,
                "amount": u.amount,
                "num_addresses": len(u.addresses),
                "locktime": pd.Timestamp(u.locktime, unit="s"),
            }
            for u in utxos
        ])

class XChainClient:
    async def get_utxos_df(
        self,
        addresses: List[str],
    ) -> pd.DataFrame:
        """Get UTXOs as DataFrame."""
        utxos = await self.get_utxos(addresses)
        return UTXO.to_dataframe(utxos)
\end{lstlisting}

\subsection{Transaction History}

\begin{lstlisting}[language=Python,caption=Transaction DataFrame]
@dataclass
class Transaction:
    tx_id: str
    chain: str
    type: str
    timestamp: int
    fee: int
    inputs: List[dict]
    outputs: List[dict]

    @classmethod
    def to_dataframe(cls, txs: List["Transaction"]) -> pd.DataFrame:
        """Convert transactions to analysis-ready DataFrame."""
        records = []
        for tx in txs:
            # Flatten inputs and outputs
            total_in = sum(i.get("amount", 0) for i in tx.inputs)
            total_out = sum(o.get("amount", 0) for o in tx.outputs)

            records.append({
                "tx_id": tx.tx_id,
                "chain": tx.chain,
                "type": tx.type,
                "timestamp": pd.Timestamp(tx.timestamp, unit="s"),
                "fee": tx.fee,
                "total_in": total_in,
                "total_out": total_out,
                "num_inputs": len(tx.inputs),
                "num_outputs": len(tx.outputs),
            })

        df = pd.DataFrame(records)
        df.set_index("timestamp", inplace=True)
        return df.sort_index()

class HistoryClient:
    async def get_address_history(
        self,
        address: str,
        start_time: pd.Timestamp | None = None,
        end_time: pd.Timestamp | None = None,
    ) -> pd.DataFrame:
        """Get transaction history for address."""
        params = {"address": address}
        if start_time:
            params["startTime"] = int(start_time.timestamp())
        if end_time:
            params["endTime"] = int(end_time.timestamp())

        txs = await self._fetch_all_pages("getAddressHistory", params)
        return Transaction.to_dataframe(txs)
\end{lstlisting}

\subsection{Validator Analytics}

\begin{lstlisting}[language=Python,caption=Validator DataFrame]
async def get_validator_stats(client: LuxClient) -> pd.DataFrame:
    """Get comprehensive validator statistics."""
    validators = await client.p_chain.get_current_validators()

    df = pd.DataFrame([
        {
            "node_id": v.node_id,
            "start_time": pd.Timestamp(v.start_time, unit="s"),
            "end_time": pd.Timestamp(v.end_time, unit="s"),
            "stake_amount": v.stake_amount / 1e9,  # Convert to LUX
            "delegation_fee": v.delegation_fee,
            "uptime": v.uptime,
            "connected": v.connected,
            "delegator_count": len(v.delegators),
            "delegated_stake": sum(d.stake for d in v.delegators) / 1e9,
        }
        for v in validators
    ])

    # Add computed columns
    df["total_stake"] = df["stake_amount"] + df["delegated_stake"]
    df["remaining_days"] = (df["end_time"] - pd.Timestamp.now()).dt.days
    df["stake_share"] = df["total_stake"] / df["total_stake"].sum()

    return df.set_index("node_id").sort_values("total_stake", ascending=False)
\end{lstlisting}

\section{Batch Operations}

\subsection{Parallel Data Fetching}

\begin{lstlisting}[language=Python,caption=Batch Fetching]
import asyncio
from typing import TypeVar, Callable, List
from itertools import islice

T = TypeVar("T")

async def batch_fetch(
    items: List[T],
    fetch_fn: Callable[[T], asyncio.Future],
    batch_size: int = 100,
    max_concurrent: int = 10,
) -> List:
    """Fetch items in parallel batches."""
    semaphore = asyncio.Semaphore(max_concurrent)

    async def fetch_with_limit(item: T):
        async with semaphore:
            return await fetch_fn(item)

    results = []
    for i in range(0, len(items), batch_size):
        batch = items[i:i + batch_size]
        batch_results = await asyncio.gather(
            *[fetch_with_limit(item) for item in batch],
            return_exceptions=True,
        )
        results.extend(batch_results)

    # Filter out exceptions
    return [r for r in results if not isinstance(r, Exception)]

# Usage
async def fetch_all_balances(
    client: LuxClient,
    addresses: List[str],
) -> pd.DataFrame:
    """Fetch balances for many addresses efficiently."""

    async def fetch_balance(addr: str) -> dict:
        balance = await client.get_balance(addr)
        return {"address": addr, "balance": balance}

    results = await batch_fetch(
        addresses,
        fetch_balance,
        batch_size=100,
        max_concurrent=20,
    )

    return pd.DataFrame(results)
\end{lstlisting}

\subsection{Historical Data Collection}

\begin{lstlisting}[language=Python,caption=Historical Collection]
from datetime import datetime, timedelta
import pyarrow as pa
import pyarrow.parquet as pq

async def collect_historical_txs(
    client: LuxClient,
    start_date: datetime,
    end_date: datetime,
    output_path: str,
) -> None:
    """Collect historical transactions to Parquet files."""

    current = start_date
    while current < end_date:
        next_day = current + timedelta(days=1)

        # Fetch day's transactions
        txs = await client.get_transactions(
            start_time=current,
            end_time=next_day,
        )

        if txs:
            # Convert to DataFrame and save
            df = Transaction.to_dataframe(txs)

            # Partition by date
            date_str = current.strftime("%Y-%m-%d")
            file_path = f"{output_path}/date={date_str}/data.parquet"

            table = pa.Table.from_pandas(df)
            pq.write_table(table, file_path, compression="snappy")

            print(f"Saved {len(df)} transactions for {date_str}")

        current = next_day

# Query collected data with DuckDB
import duckdb

def query_historical_data(data_path: str, query: str) -> pd.DataFrame:
    """Query Parquet files with SQL."""
    conn = duckdb.connect()
    return conn.execute(f"""
        SELECT * FROM read_parquet('{data_path}/**/*.parquet')
        WHERE {query}
    """).fetchdf()

# Example: Get daily transaction volumes
volumes = query_historical_data(
    "/data/lux/transactions",
    """
    SELECT
        DATE_TRUNC('day', timestamp) as date,
        COUNT(*) as tx_count,
        SUM(total_in) / 1e9 as volume_lux
    FROM transactions
    GROUP BY 1
    ORDER BY 1
    """
)
\end{lstlisting}

\section{Streaming Data}

\subsection{WebSocket Subscriptions}

\begin{lstlisting}[language=Python,caption=WebSocket Streaming]
import asyncio
from collections import deque
from typing import Callable

class StreamProcessor:
    """Process blockchain events in real-time."""

    def __init__(self, client: LuxClient, buffer_size: int = 1000):
        self.client = client
        self.buffer = deque(maxlen=buffer_size)
        self.handlers: List[Callable] = []

    def on_event(self, handler: Callable):
        """Register event handler."""
        self.handlers.append(handler)
        return handler

    async def start(self, chain: str = "X"):
        """Start processing stream."""
        async for block in self.client.stream_blocks(chain):
            self.buffer.append(block)

            for handler in self.handlers:
                try:
                    await handler(block)
                except Exception as e:
                    print(f"Handler error: {e}")

    def recent_blocks(self, n: int = 100) -> List[Block]:
        """Get recent blocks from buffer."""
        return list(self.buffer)[-n:]

# Usage
processor = StreamProcessor(client)

@processor.on_event
async def log_block(block: Block):
    print(f"New block: {block.id} with {len(block.txs)} txs")

@processor.on_event
async def update_metrics(block: Block):
    # Update Prometheus metrics, etc.
    pass

await processor.start("X")
\end{lstlisting}

\subsection{Real-Time Analytics}

\begin{lstlisting}[language=Python,caption=Streaming Analytics]
from collections import defaultdict
import time

class RealTimeAnalytics:
    """Compute rolling statistics from stream."""

    def __init__(self, window_seconds: int = 300):
        self.window = window_seconds
        self.tx_counts: deque = deque()
        self.volumes: deque = deque()
        self.active_addresses: defaultdict = defaultdict(int)

    def process_block(self, block: Block):
        """Process block and update statistics."""
        now = time.time()

        # Remove old data outside window
        while self.tx_counts and self.tx_counts[0][0] < now - self.window:
            self.tx_counts.popleft()
        while self.volumes and self.volumes[0][0] < now - self.window:
            self.volumes.popleft()

        # Add new data
        for tx in block.transactions:
            self.tx_counts.append((now, 1))
            self.volumes.append((now, tx.total_value))

            for addr in tx.all_addresses:
                self.active_addresses[addr] = now

        # Clean old addresses
        self.active_addresses = {
            addr: ts
            for addr, ts in self.active_addresses.items()
            if ts > now - self.window
        }

    @property
    def tps(self) -> float:
        """Transactions per second in window."""
        if not self.tx_counts:
            return 0.0
        return len(self.tx_counts) / self.window

    @property
    def volume_per_minute(self) -> float:
        """Volume per minute in window."""
        if not self.volumes:
            return 0.0
        total = sum(v for _, v in self.volumes)
        return (total / self.window) * 60

    @property
    def active_address_count(self) -> int:
        """Unique addresses active in window."""
        return len(self.active_addresses)

    def to_dict(self) -> dict:
        """Export current metrics."""
        return {
            "tps": self.tps,
            "volume_per_minute": self.volume_per_minute,
            "active_addresses": self.active_address_count,
            "window_seconds": self.window,
        }
\end{lstlisting}

\section{Machine Learning Integration}

\subsection{Feature Engineering}

\begin{lstlisting}[language=Python,caption=Feature Engineering]
import numpy as np
from sklearn.preprocessing import StandardScaler

class TransactionFeatures:
    """Extract ML features from transactions."""

    def __init__(self):
        self.scaler = StandardScaler()

    def extract_features(self, df: pd.DataFrame) -> np.ndarray:
        """Extract numerical features from transaction DataFrame."""
        features = pd.DataFrame()

        # Time-based features
        features["hour"] = df.index.hour
        features["day_of_week"] = df.index.dayofweek
        features["is_weekend"] = (df.index.dayofweek >= 5).astype(int)

        # Value features
        features["log_value"] = np.log1p(df["total_in"])
        features["fee_ratio"] = df["fee"] / (df["total_in"] + 1)

        # Complexity features
        features["input_count"] = df["num_inputs"]
        features["output_count"] = df["num_outputs"]
        features["io_ratio"] = df["num_inputs"] / (df["num_outputs"] + 1)

        # Rolling statistics (requires sorted index)
        for window in [10, 100, 1000]:
            features[f"value_ma_{window}"] = (
                df["total_in"].rolling(window, min_periods=1).mean()
            )
            features[f"value_std_{window}"] = (
                df["total_in"].rolling(window, min_periods=1).std().fillna(0)
            )
            features[f"tx_count_{window}"] = (
                df["tx_id"].rolling(window, min_periods=1).count()
            )

        return features.values

    def fit_transform(self, df: pd.DataFrame) -> np.ndarray:
        """Fit scaler and transform features."""
        features = self.extract_features(df)
        return self.scaler.fit_transform(features)

    def transform(self, df: pd.DataFrame) -> np.ndarray:
        """Transform features using fitted scaler."""
        features = self.extract_features(df)
        return self.scaler.transform(features)
\end{lstlisting}

\subsection{Anomaly Detection}

\begin{lstlisting}[language=Python,caption=Anomaly Detection]
from sklearn.ensemble import IsolationForest
import joblib

class TransactionAnomalyDetector:
    """Detect anomalous transactions using Isolation Forest."""

    def __init__(self, contamination: float = 0.01):
        self.feature_extractor = TransactionFeatures()
        self.model = IsolationForest(
            contamination=contamination,
            random_state=42,
            n_jobs=-1,
        )

    def fit(self, df: pd.DataFrame):
        """Train on historical transactions."""
        features = self.feature_extractor.fit_transform(df)
        self.model.fit(features)
        return self

    def predict(self, df: pd.DataFrame) -> pd.Series:
        """Predict anomaly scores (-1 = anomaly, 1 = normal)."""
        features = self.feature_extractor.transform(df)
        predictions = self.model.predict(features)
        return pd.Series(predictions, index=df.index, name="anomaly")

    def score(self, df: pd.DataFrame) -> pd.Series:
        """Get anomaly scores (lower = more anomalous)."""
        features = self.feature_extractor.transform(df)
        scores = self.model.score_samples(features)
        return pd.Series(scores, index=df.index, name="anomaly_score")

    def save(self, path: str):
        """Save model to disk."""
        joblib.dump({
            "feature_extractor": self.feature_extractor,
            "model": self.model,
        }, path)

    @classmethod
    def load(cls, path: str) -> "TransactionAnomalyDetector":
        """Load model from disk."""
        data = joblib.load(path)
        detector = cls()
        detector.feature_extractor = data["feature_extractor"]
        detector.model = data["model"]
        return detector

# Training pipeline
async def train_anomaly_detector(
    client: LuxClient,
    days: int = 30,
) -> TransactionAnomalyDetector:
    """Train detector on recent transaction history."""
    end_time = pd.Timestamp.now()
    start_time = end_time - pd.Timedelta(days=days)

    # Fetch historical data
    df = await client.get_transaction_history(
        start_time=start_time,
        end_time=end_time,
    )

    # Train model
    detector = TransactionAnomalyDetector()
    detector.fit(df)

    return detector
\end{lstlisting}

\subsection{Trading Signals}

\begin{lstlisting}[language=Python,caption=Trading Signal Generation]
from dataclasses import dataclass
from enum import Enum

class SignalType(Enum):
    BUY = "buy"
    SELL = "sell"
    HOLD = "hold"

@dataclass
class TradingSignal:
    timestamp: pd.Timestamp
    signal_type: SignalType
    confidence: float
    features: dict

class OnChainSignalGenerator:
    """Generate trading signals from on-chain data."""

    def __init__(self, client: LuxClient):
        self.client = client
        self.analytics = RealTimeAnalytics(window_seconds=300)
        self.history: deque = deque(maxlen=1000)

    async def process_block(self, block: Block) -> TradingSignal | None:
        """Process block and generate signal if conditions met."""
        self.analytics.process_block(block)
        self.history.append(self.analytics.to_dict())

        if len(self.history) < 100:
            return None  # Need more data

        # Extract features
        tps_history = [h["tps"] for h in self.history]
        volume_history = [h["volume_per_minute"] for h in self.history]

        # Simple momentum strategy
        tps_ma_short = np.mean(tps_history[-10:])
        tps_ma_long = np.mean(tps_history[-100:])
        volume_ma_short = np.mean(volume_history[-10:])
        volume_ma_long = np.mean(volume_history[-100:])

        # Generate signal
        tps_momentum = tps_ma_short / (tps_ma_long + 1e-9)
        volume_momentum = volume_ma_short / (volume_ma_long + 1e-9)

        features = {
            "tps_momentum": tps_momentum,
            "volume_momentum": volume_momentum,
            "active_addresses": self.analytics.active_address_count,
        }

        if tps_momentum > 1.2 and volume_momentum > 1.3:
            return TradingSignal(
                timestamp=pd.Timestamp.now(),
                signal_type=SignalType.BUY,
                confidence=min(tps_momentum * volume_momentum / 2, 1.0),
                features=features,
            )
        elif tps_momentum < 0.8 and volume_momentum < 0.7:
            return TradingSignal(
                timestamp=pd.Timestamp.now(),
                signal_type=SignalType.SELL,
                confidence=min((2 - tps_momentum) * (2 - volume_momentum) / 2, 1.0),
                features=features,
            )

        return None

    async def run(self):
        """Run signal generation loop."""
        async for block in self.client.stream_blocks():
            signal = await self.process_block(block)
            if signal:
                yield signal
\end{lstlisting}

\section{Caching and Persistence}

\subsection{Local Cache}

\begin{lstlisting}[language=Python,caption=Caching Layer]
import sqlite3
import pickle
from functools import wraps
from typing import Callable
import hashlib

class BlockchainCache:
    """SQLite-based cache for blockchain data."""

    def __init__(self, db_path: str = "lux_cache.db"):
        self.conn = sqlite3.connect(db_path)
        self._init_tables()

    def _init_tables(self):
        self.conn.execute("""
            CREATE TABLE IF NOT EXISTS cache (
                key TEXT PRIMARY KEY,
                value BLOB,
                expires_at REAL
            )
        """)
        self.conn.commit()

    def get(self, key: str) -> any | None:
        """Get cached value."""
        cursor = self.conn.execute(
            "SELECT value, expires_at FROM cache WHERE key = ?",
            (key,)
        )
        row = cursor.fetchone()

        if row is None:
            return None

        value, expires_at = row
        if expires_at and time.time() > expires_at:
            self.delete(key)
            return None

        return pickle.loads(value)

    def set(self, key: str, value: any, ttl: int | None = None):
        """Set cached value with optional TTL."""
        expires_at = time.time() + ttl if ttl else None
        self.conn.execute(
            "INSERT OR REPLACE INTO cache (key, value, expires_at) VALUES (?, ?, ?)",
            (key, pickle.dumps(value), expires_at)
        )
        self.conn.commit()

    def delete(self, key: str):
        """Delete cached value."""
        self.conn.execute("DELETE FROM cache WHERE key = ?", (key,))
        self.conn.commit()

def cached(ttl: int = 300):
    """Decorator for caching async function results."""
    def decorator(func: Callable):
        @wraps(func)
        async def wrapper(self, *args, **kwargs):
            # Generate cache key
            key_data = f"{func.__name__}:{args}:{kwargs}"
            key = hashlib.sha256(key_data.encode()).hexdigest()

            # Check cache
            cached_value = self._cache.get(key)
            if cached_value is not None:
                return cached_value

            # Call function
            result = await func(self, *args, **kwargs)

            # Store in cache
            self._cache.set(key, result, ttl)

            return result
        return wrapper
    return decorator

class CachedClient(LuxClient):
    """Client with automatic caching."""

    def __init__(self, *args, cache_path: str = "lux_cache.db", **kwargs):
        super().__init__(*args, **kwargs)
        self._cache = BlockchainCache(cache_path)

    @cached(ttl=60)
    async def get_balance(self, address: str, asset_id: str = "LUX") -> int:
        return await super().get_balance(address, asset_id)

    @cached(ttl=300)
    async def get_validator_stats(self) -> pd.DataFrame:
        return await get_validator_stats(self)
\end{lstlisting}

\section{Benchmarks}

\subsection{Data Fetching Performance}

Measured fetching from local node, Python 3.11, M1 MacBook Pro:

\begin{table}[h]
\centering
\caption{Fetch Performance (1000 requests)}
\begin{tabular}{lrrr}
\toprule
\textbf{Operation} & \textbf{Sequential} & \textbf{Parallel (20)} & \textbf{Speedup} \\
\midrule
get\_balance & 12.3 s & 0.82 s & 15x \\
get\_utxos & 45.2 s & 3.1 s & 14.6x \\
get\_tx\_status & 8.5 s & 0.61 s & 13.9x \\
\bottomrule
\end{tabular}
\end{table}

\subsection{DataFrame Operations}

\begin{table}[h]
\centering
\caption{DataFrame Conversion Performance}
\begin{tabular}{lrr}
\toprule
\textbf{Records} & \textbf{Time} & \textbf{Memory} \\
\midrule
1,000 & 2.1 ms & 0.4 MB \\
10,000 & 18 ms & 3.8 MB \\
100,000 & 180 ms & 38 MB \\
1,000,000 & 1.9 s & 380 MB \\
\bottomrule
\end{tabular}
\end{table}

\subsection{Rust vs Pure Python}

\begin{table}[h]
\centering
\caption{Rust Core vs Pure Python}
\begin{tabular}{lrrr}
\toprule
\textbf{Operation} & \textbf{Rust} & \textbf{Python} & \textbf{Speedup} \\
\midrule
Sign transaction & 0.18 ms & 8.5 ms & 47x \\
Parse UTXO (1000) & 0.4 ms & 45 ms & 112x \\
Address encoding & 0.02 ms & 1.2 ms & 60x \\
\bottomrule
\end{tabular}
\end{table}

\subsection{ML Pipeline}

\begin{table}[h]
\centering
\caption{ML Pipeline Performance (100K transactions)}
\begin{tabular}{lr}
\toprule
\textbf{Step} & \textbf{Time} \\
\midrule
Feature extraction & 2.3 s \\
Model training (IsolationForest) & 4.1 s \\
Prediction (batch) & 0.8 s \\
Prediction (single) & 0.3 ms \\
\bottomrule
\end{tabular}
\end{table}

\section{Usage Examples}

\subsection{Basic Analytics}

\begin{lstlisting}[language=Python,caption=Basic Usage]
import asyncio
from luxsdk import LuxClient

async def main():
    async with LuxClient("https://api.lux.network") as client:
        # Get validator statistics
        validators = await client.p_chain.get_validator_stats()
        print(f"Total validators: {len(validators)}")
        print(f"Total staked: {validators['total_stake'].sum():.2f} LUX")

        # Analyze top validators
        top_10 = validators.head(10)
        print("\nTop 10 validators by stake:")
        print(top_10[["stake_amount", "delegated_stake", "uptime"]])

        # Get address history
        history = await client.get_address_history(
            "X-lux1...",
            start_time=pd.Timestamp("2023-01-01"),
        )

        # Daily transaction volume
        daily = history.resample("D").agg({
            "total_in": "sum",
            "tx_id": "count",
        })
        print("\nDaily statistics:")
        print(daily.describe())

asyncio.run(main())
\end{lstlisting}

\subsection{Real-Time Monitoring}

\begin{lstlisting}[language=Python,caption=Real-Time Monitoring]
import asyncio
from luxsdk import LuxClient
from luxsdk.ml import TransactionAnomalyDetector

async def monitor_anomalies():
    """Monitor blockchain for anomalous transactions."""
    async with LuxClient("https://api.lux.network") as client:
        # Load pre-trained detector
        detector = TransactionAnomalyDetector.load("models/detector.joblib")

        async for block in client.stream_blocks():
            if not block.transactions:
                continue

            # Convert to DataFrame
            df = Transaction.to_dataframe(block.transactions)

            # Score transactions
            scores = detector.score(df)

            # Alert on anomalies
            anomalies = df[scores < -0.5]
            for _, tx in anomalies.iterrows():
                print(f"ANOMALY: {tx.name} - score: {scores[tx.name]:.3f}")
                # Send alert, log, etc.

asyncio.run(monitor_anomalies())
\end{lstlisting}

\section{Related Work}

The web3.py library~\cite{web3py} established patterns for Python blockchain SDKs. The pandas library~\cite{pandas} provides the foundation for data analysis. Scikit-learn~\cite{sklearn} enables accessible ML pipelines. Our SDK integrates these tools specifically for Lux blockchain analytics.

\section{Conclusion}

The Lux Python SDK provides a comprehensive toolkit for blockchain data science and ML applications. Key contributions:

\begin{itemize}
    \item Native pandas integration for exploratory analysis
    \item Async batch operations with 15x throughput improvement
    \item Real-time streaming for live analytics
    \item ML pipeline integration with feature engineering
    \item Rust core for 50-100x speedup on crypto operations
    \item SQLite caching for efficient repeated queries
\end{itemize}

The SDK bridges blockchain infrastructure and the Python data science ecosystem, enabling new applications in quantitative analysis, fraud detection, and automated trading. Future work includes GPU-accelerated feature extraction, integration with PyTorch for deep learning, and distributed processing with Dask.

\begin{thebibliography}{9}

\bibitem{web3py}
Ethereum Foundation.
\textit{web3.py: A Python library for interacting with Ethereum}.
\url{https://web3py.readthedocs.io}, 2016.

\bibitem{pandas}
McKinney, W.
\textit{pandas: Powerful data structures for data analysis}.
\url{https://pandas.pydata.org}, 2008.

\bibitem{sklearn}
Pedregosa, F., et al.
\textit{Scikit-learn: Machine Learning in Python}.
Journal of Machine Learning Research, 2011.

\bibitem{pyo3}
PyO3 Contributors.
\textit{PyO3: Rust bindings for Python}.
\url{https://pyo3.rs}, 2017.

\bibitem{asyncio}
Python Software Foundation.
\textit{asyncio: Asynchronous I/O}.
Python Standard Library, 2014.

\end{thebibliography}

\end{document}
